\documentclass[11pt]{article}
\usepackage[margin=1in]{geometry}
\usepackage{authblk}
\usepackage{amsmath, amssymb}
\usepackage{booktabs}
\usepackage{enumitem}
\usepackage[hidelinks]{hyperref}
\usepackage{xcolor}
\definecolor{flag}{HTML}{8b0000}
\newcommand{\dkl}{\Delta KL}

\title{Request for criticism:\\
have we established that trademark filings measure\\
anything --- innovation, business formation, or otherwise?}
\author{Eric Silver}
\affil{}
\date{May 2026}

\begin{document}
\maketitle

\noindent
\textbf{What I am asking you to do.} Attack the construct validity of
this body of work. I am not asking whether the regressions are coded
correctly (assume they are; the code is reproducible) or whether the
prose is polished. I am asking the one question that, if answered
badly, makes everything downstream worthless: \emph{does a trademark
filing, and the $\dkl$ statistic I compute on its goods/services
text, measure anything real --- innovation, business formation, or
even meaningful commercial activity --- or am I measuring the
drafting habits of trademark attorneys and calling it economics?}

I have written this to make your job easy. Below I state the claims,
then I steelman the case that the whole edifice is measuring noise,
then I lay out the thin evidence I have for validity, then I offer
four reframings from strong to weak and ask you to tell me which (if
any) survives. Be as brutal as the strongest version of the
objection warrants. I would rather kill the ``innovation'' framing
now than have a referee or a serious reader kill it later.

\section{The claims, briefly}

There are two bodies of work sharing one measurement apparatus.

\paragraph{(A) The $\dkl$ papers (\texttt{main.pdf}, \texttt{short.pdf}).}
For each of $10{,}060{,}073$ ``clean'' US trademark filings
(1990--2021 cohorts), I compute a Kullback--Leibler surprise
difference $\dkl = KL_{\text{pros}} - KL_{\text{retr}}$ on the
goods/services description text, against per-day exponentially
decay-weighted reference distributions of the same NICE class
(half-life $H = 2$y). $\dkl > 0$ when a filing's vocabulary is
unusual against the recent past and close to the near future --- the
directional resonance signal of Murdock, Allen, \& DeDeo (2017),
applied for the first time to commercially-incentivised text. The
papers claim: $\dkl$ is an innovation measure; it predicts
post-registration survival ($U$-shaped) and financial outcomes; it
is statistically independent of patent counts ($r = +0.01$ on
$174{,}569$ firm-years) and of comparable correlation magnitude;
the innovation tail aligns with the empirical content of optimal
distinctiveness and the harvest tail is an unexplained surprise.

\paragraph{(B) The dynamism / ethnic-cluster work
(\texttt{dynamism.pdf}, \texttt{ethnic\_clusters\_note.pdf}).}
Treating a firm's first-ever filing (``debut'') as a proxy for
business formation, I show trademark debuts cover $8$--$18\%$ of
Census High-Propensity Business Applications; I impute applicant
race/ethnicity from surnames; I document ethnic-cluster industries,
a Black trademark-rise / patent-fall divergence, and the collapse of
an apparent ``patent guild'' signal once foreign inventors are
stripped.

Everything in (B) inherits the validity of the measurement in (A),
plus its own surname-imputation problems. So the load-bearing
question is shared.

\section{Steelman: the case that this measures nothing real}

Here is the argument I most want you to either complete or refute. I
find it uncomfortably strong. \textbf{[moderate]} that at least the
first three links are serious and not fully answered by the current
drafts.

\begin{enumerate}[leftmargin=*]
\item \textbf{A trademark is a brand, not an invention, and the
goods/services text is a legal instrument, not a description of
novelty.} The g/s description exists to define the \emph{scope of
protection}, drafted by an attorney (or selected from the USPTO's
pre-approved \emph{ID Manual}) to be as broad as the examiner will
allow. Two firms selling physically identical products can write
wildly different descriptions; a firm with a genuinely novel product
can file boilerplate. So $\dkl$ measures the novelty of the
\emph{prose}, and the prose is several steps removed from the
product. Why should I believe the prose tracks the innovation?

\item \textbf{Vocabulary novelty has at least four generators, and I
have separated none of them.} A high-$\dkl$ filing uses words that
are unusual-then-common. That could be (a) a genuine new product
category (crypto, vape, NFT); (b) legal-drafting fashion; (c)
attorneys copying each other's successful protective language; (d)
the USPTO updating its ID Manual, which mechanically shifts the
reference vocabulary for everyone. The cannabis/crypto examples I
parade are consistent with (a), but they are also exactly the cases
where new ID-Manual entries appeared. I have not shown the signal
is (a) rather than (d).

\item \textbf{The resonance interpretation was validated on a
different species of text.} Murdock, Allen, \& DeDeo (2017,
\emph{Cognition}) developed the framework on Darwin's private
reading notebooks; Barron, Huang, Spang, \& DeDeo (2018,
\emph{PNAS}) on French Revolution parliamentary debates --- in both,
the text is a window into the author's cognitive or rhetorical
state, authored by the person whose novelty is being measured. A
trademark
g/s description is drafted for protection scope by a third-party
agent. Importing the ``resonance = the field moved toward this
filing'' interpretation assumes the signal means the same thing in
a corpus where the author is a lawyer optimising a legal outcome. I
assert the analogy; I do not re-validate it.

\item \textbf{Selection: I only ever observe the $\sim 8$--$18\%$ of
business starts that trademark.} Coverage against Census
High-Propensity Business Applications is $8\%$ (2004--08) rising to
$18\%$ (2019--24). The trademark-filing population is selected on
branding intensity, consumer-facing orientation, and capitalisation.
Every sentence I write ``about firms'' is really about the
trademark-filing subset, and that subset is exactly the set of firms
for which branding-vocabulary is salient --- which is precisely the
variable $\dkl$ measures. The selection is on the regressor.

\item \textbf{``Survival'' is trademark maintenance, not firm
survival.} My headline outcome --- 5-year survival --- is the owner
paying the \S\,8 declaration-of-use fee. A live firm can abandon a
mark; a dead firm's mark can sit registered until renewal lapses.
The outcome is a portfolio-management decision correlated with, but
not identical to, firm viability. The $U$-shape on $\dkl$ could be a
$U$-shape in \emph{which marks firms bother to maintain}, not in
which firms survive.

\item \textbf{The financial validation is small and one of the
numbers is implausible.} The patent $R^2$ decomposition is on
$8{,}118$ firm-years; the gross-margin effect is $+3.2$pp per
$\sigma$ (SE $0.6$); these are real but modest. The claimed
$+28$ percentage-point 4-year excess return per $\sigma$ of debut
$\dkl$ is, on its face, \textbf{[low]} confidence and probably too
large to believe --- a $\sigma$ move in a text statistic should not
move 4-year returns by $28$pp; I suspect a specification or
sample-composition artefact and want it stress-tested.

\item \textbf{Independence from patents cuts both ways.} I sell
$r = +0.01$ with patent counts as evidence that $\dkl$ is ``a
distinct innovation channel.'' The skeptic's reading is the
opposite: $\dkl$ is uncorrelated with the one measure we
\emph{know} tracks invention because $\dkl$ is not measuring
invention at all --- it is measuring marketing/branding activity,
which is orthogonal to R\&D by construction. Independence is
necessary for ``distinct channel'' but equally consistent with
``not an innovation measure.''

\item \textbf{The ethnic / dynamism layer multiplies the problem.}
Surname imputation is probabilistic, restricted to the $\sim 32\%$
of filers parsing as individuals, static to the 2010 Census, and
weakest exactly for Black surnames (Smith/Johnson/Williams load on
both White and Black). On top of that it inherits the whole
proxy-validity question: if trademark debut is not business
formation, the coverage ratios and ethnic-composition trends
describe ``who files trademarks,'' not ``who starts businesses.''
\end{enumerate}

\section{What I actually have for validity (honest accounting)}

I am not claiming the steelman wins. Here is the evidence on the
other side, with its limits stated.

\begin{itemize}[leftmargin=*]
\item \textbf{Patent independence + comparable magnitude.}
$r = +0.01$ with patent counts ($174{,}569$ firm-years); incremental
4-year-return $R^2$ of $+0.0046$ for $\dkl$ vs $+0.0024$ for
$\log(1+\text{patents})$, roughly additive. \emph{Limit:} small
$R^2$ absolute; comparable-to-patents is a low bar if patents
themselves explain little return variance in this panel.
\item \textbf{Survival is $U$-shaped under multiple
operationalisations and in $32/44$ NICE classes.} \emph{Limit:}
all of them are maintenance-based; the robustness is across
definitions of maintenance, not across independent firm-outcome
data.
\item \textbf{Cluster-viability mediation absorbs $84\%$ of the
quadratic.} This is the one result that looks mechanism-like rather
than mechanical. \emph{Limit:} the mediator is itself built from
trademark data (cosine-neighbour survival), so it does not break out
of the corpus.
\item \textbf{Coverage brackets Dinlersoz et al.\ (2018).} Our
$8$--$18\%$ spans their restricted-LBD $10$--$15\%$. \emph{Limit:}
agreement on the \emph{coverage fraction} is not validation that the
covered filings \emph{mean} anything; it just says we count debut
filers the way Census counts them.
\item \textbf{Face validity.} Canonical winners skew positive
$\dkl$; flux-neutral tokens (\texttt{organisational},
\texttt{conformity}) die. \emph{Limit:} selected ex-post; not a
test.
\end{itemize}

\noindent
My own read \textbf{[moderate]}: I can defend ``$\dkl$ is a
nontrivial, patent-orthogonal signal that predicts trademark-level
outcomes and correlates weakly with firm financials.'' I cannot
currently defend ``$\dkl$ is an innovation measure'' against a
determined skeptic, because every outcome I validate against is
either inside the trademark corpus or a small noisy financial panel.

\section{Four reframings --- which survives?}

I want you to tell me which of these I can defend, and whether the
strong claim is salvageable with additional work or should be
abandoned.

\begin{description}[leftmargin=*]
\item[(I) Strong --- ``$\dkl$ is an innovation measure.''] The
current framing. Requires that vocabulary novelty in g/s prose
tracks product/market innovation. My confidence this survives
review as stated: \textbf{[low]}.
\item[(II) Medium --- ``$\dkl$ is a leading indicator of commercial
category emergence.''] Claims the corpus-level signal (which
\emph{vocabularies} are forming and being adopted) is real and
useful, while staying agnostic about whether any individual filer is
``innovative.'' The cannabis/crypto/NFT waves are the evidence; the
firm-level survival/returns become secondary. Confidence:
\textbf{[moderate]}.
\item[(III) Weak/descriptive --- ``$\dkl$ measures the linguistic
position of a filing relative to its NICE class's vocabulary
trajectory.''] Purely descriptive. Every downstream correlation
(survival, returns, ethnic clusters) becomes an empirical regularity
to be explained, not a validation of an innovation construct. This
is almost certainly defensible \textbf{[high]} but gives up the
headline.
\item[(IV) Activity proxy --- ``a trademark debut is formal
commercial activity with branding intent.''] For the dynamism work
specifically: not innovation, not business formation \emph{per se},
but a measurable, name-disclosed, near-real-time signal of
branding-intent commercial activity, explicitly the $8$--$18\%$
branding-intensive slice of formation. Confidence: \textbf{[moderate]}.
\end{description}

\noindent
A plausible restructure: lead with (III) as the measurement
contribution, establish (II) at the corpus level with the
category-emergence evidence, and present the firm-level
innovation/financial results as \emph{suggestive regularities under
(III)} rather than as proof of (I). The dynamism work then stands on
(IV) and does not need (I) at all. Tell me if this is cowardice or
rigor.

\section{Specific questions}

\begin{enumerate}[leftmargin=*]
\item Is link 1 (prose $\neq$ product) fatal, or is there a
defensible argument that attorney-drafted scope language still
carries firm-level signal because the firm chooses what to claim?
\item What is the cheapest decisive test that would separate
``genuine category emergence'' from ``ID-Manual / drafting-fashion
drift''? (My best idea: hold out the ID-Manual update history and
show $\dkl$ predicts adoption \emph{before} the manual codifies the
term. I have not done this.)
\item Is there any external, non-trademark, non-financial outcome I
could validate survival against --- business registration death,
UCC-lien termination, Form-D withdrawal, domain expiry --- that
would break the maintenance-vs-survival circularity?
\item Does the $+28$pp/$\sigma$ return number kill my credibility on
sight? Should I drop it pending re-estimation?
\item For the ethnic work: is surname imputation usable at all given
the Black-surname cross-loading, or should I restrict to the
Hispanic and API analyses where the Census classifier is
discriminating and drop Black claims to a flagged-caveat appendix?
\item Is reframing (III)+(II)+(IV) the right retreat, or is it
giving up a defensible (I) too cheaply?
\end{enumerate}

\section{Artifacts and reproduction}

All on branch \texttt{business-dynamism} of the project repo.
\texttt{main.pdf} (long $\dkl$ paper, 223pp), \texttt{short.pdf}
(7pp), \texttt{dynamism.pdf} (13pp), \texttt{ethnic\_clusters\_note.pdf}
(21pp). Pipelines: \texttt{Makefile} and \texttt{Makefile.dynamism}.
Every figure regenerates from \texttt{scripts/} against the
per-NICE-class parquets; the surname work uses the 2010 Census
surname file and a full re-parse of the USPTO XML for applicant
country (\texttt{owner\_address.parquet}, $13{,}993{,}768$ records).

\section*{References}
\begin{itemize}[leftmargin=*]\setlength\itemsep{2pt}
\item Akcigit, U., \& Ates, S.~T. (2021). Ten facts on declining
business dynamism and lessons from endogenous growth theory.
\emph{American Economic Journal: Macroeconomics}, 13(1), 257--298.
\item Barron, A.~T.~J., Huang, J., Spang, R.~L., \& DeDeo, S. (2018).
Individuals, institutions, and innovation in the debates of the
French Revolution. \emph{PNAS}, 115(18), 4607--4612.
\item Dinlersoz, E., Goldschlag, N., Myers, A., \& Zolas, N. (2018).
An anatomy of US firms seeking trademark registration. Center for
Economic Studies WP 18--46, US Census Bureau.
\item Decker, R.~A., Haltiwanger, J., Jarmin, R.~S., \& Miranda, J.
(2014). The role of entrepreneurship in US job creation and economic
dynamism. \emph{Journal of Economic Perspectives}, 28(3), 3--24.
\item Murdock, J., Allen, C., \& DeDeo, S. (2017). Exploration and
exploitation of Victorian science in Darwin's reading notebooks.
\emph{Cognition}, 159, 117--126.
\item Word, D.~L., Coleman, C.~D., Nunziata, R., \& Kominski, R.
(2008). Demographic aspects of surnames from Census 2000. US Census
Bureau working paper.
\end{itemize}

\end{document}
